% ১১. Form
\section{Form}
Form হলো ডেটাবেজে ডেটা ইনপুট, দেখা ও সম্পাদনার জন্য ব্যবহারবান্ধব ইন্টারফেস। সরাসরি টেবিলে কাজ করলে ভুল হওয়ার সম্ভাবনা বেশি, তাই ব্যবহারকারীর জন্য ফর্ম তৈরি করা হয়।

\begin{definitionbox}{Form}
ডেটাবেজের টেবিল বা কুয়েরির ডেটা সহজভাবে এন্ট্রি, প্রদর্শন ও সম্পাদনার জন্য তৈরি ইন্টারফেসকে Form বলে।
\end{definitionbox}

\begin{keypointbox}{ফর্মের উপাদান}
Label, Text box, Combo box, List box, Check box, Radio button, Command button, Date picker ইত্যাদি কন্ট্রোল ফর্মে ব্যবহার করা যায়।
\end{keypointbox}

\begin{examplebox}{চিত্রের স্থানধারক}
Student Entry Form: StudentID, Name, GroupName, BirthDate, Mobile এবং Save/Search/Delete বোতামসহ একটি ফর্মের নকশা।
\end{examplebox}

\begin{boardbox}{ব্যবহারিক গুরুত্ব}
MS Access-এ টেবিলের উপর ভিত্তি করে Form Wizard বা Design View ব্যবহার করে ফর্ম তৈরি করা যায়। ফর্মের মাধ্যমে ডেটা এন্ট্রি করলে ব্যবহারকারীকে পুরো টেবিলের কাঠামো জানতে হয় না।
\end{boardbox}
